% Options for packages loaded elsewhere
\PassOptionsToPackage{unicode}{hyperref}
\PassOptionsToPackage{hyphens}{url}
\documentclass[
  12pt,
]{ctexart}
\usepackage{xcolor}
\usepackage{amsmath,amssymb}
\setcounter{secnumdepth}{-\maxdimen} % remove section numbering
\usepackage{iftex}
\ifPDFTeX
  \usepackage[T1]{fontenc}
  \usepackage[utf8]{inputenc}
  \usepackage{textcomp} % provide euro and other symbols
\else % if luatex or xetex
  \usepackage{unicode-math} % this also loads fontspec
  \defaultfontfeatures{Scale=MatchLowercase}
  \defaultfontfeatures[\rmfamily]{Ligatures=TeX,Scale=1}
\fi
\usepackage{lmodern}
\ifPDFTeX\else
  % xetex/luatex font selection
\fi
% Use upquote if available, for straight quotes in verbatim environments
\IfFileExists{upquote.sty}{\usepackage{upquote}}{}
\IfFileExists{microtype.sty}{% use microtype if available
  \usepackage[]{microtype}
  \UseMicrotypeSet[protrusion]{basicmath} % disable protrusion for tt fonts
}{}
\makeatletter
\@ifundefined{KOMAClassName}{% if non-KOMA class
  \IfFileExists{parskip.sty}{%
    \usepackage{parskip}
  }{% else
    \setlength{\parindent}{0pt}
    \setlength{\parskip}{6pt plus 2pt minus 1pt}}
}{% if KOMA class
  \KOMAoptions{parskip=half}}
\makeatother
\usepackage{longtable,booktabs,array}
\usepackage{calc} % for calculating minipage widths
% Correct order of tables after \paragraph or \subparagraph
\usepackage{etoolbox}
\makeatletter
\patchcmd\longtable{\par}{\if@noskipsec\mbox{}\fi\par}{}{}
\makeatother
% Allow footnotes in longtable head/foot
\IfFileExists{footnotehyper.sty}{\usepackage{footnotehyper}}{\usepackage{footnote}}
\makesavenoteenv{longtable}
\setlength{\emergencystretch}{3em} % prevent overfull lines
\providecommand{\tightlist}{%
  \setlength{\itemsep}{0pt}\setlength{\parskip}{0pt}}
\usepackage{bookmark}
\IfFileExists{xurl.sty}{\usepackage{xurl}}{} % add URL line breaks if available
\urlstyle{same}
\hypersetup{
  hidelinks,
  pdfcreator={LaTeX via pandoc}}

\author{SCX}
\date{2026}

\begin{document}

\section{Final Mathematical Verification
Report}\label{final-mathematical-verification-report}

\textbf{Reviewer}: Pure mathematician (formal verification of
mathematical claims only) \textbf{Date}: 2026-06-28 \textbf{Scope}:
Theorems 1-5, supporting lemmas, and numerical test cases \textbf{Files
reviewed}: -
\texttt{G:\textbackslash{}Xiaogan\_Supercomputing\_data\textbackslash{}SCX\textbackslash{}theory\textbackslash{}theorems\textbackslash{}01\_noise\_detection\_guarantee.md}
-
\texttt{G:\textbackslash{}Xiaogan\_Supercomputing\_data\textbackslash{}SCX\textbackslash{}theory\textbackslash{}theorems\textbackslash{}02\_weak\_feature\_failure.md}
-
\texttt{G:\textbackslash{}Xiaogan\_Supercomputing\_data\textbackslash{}SCX\textbackslash{}theory\textbackslash{}theorems\textbackslash{}03\_unidentifiability\_theorem.md}
-
\texttt{G:\textbackslash{}Xiaogan\_Supercomputing\_data\textbackslash{}SCX\textbackslash{}theory\textbackslash{}explorations\textbackslash{}exact\_constant\_minimax.md}
-
\texttt{G:\textbackslash{}Xiaogan\_Supercomputing\_data\textbackslash{}SCX\textbackslash{}theory\textbackslash{}explorations\textbackslash{}lemma\_AB\_bahadur\_rao\_f1.md}
-
\texttt{G:\textbackslash{}Xiaogan\_Supercomputing\_data\textbackslash{}SCX\textbackslash{}theory\textbackslash{}explorations\textbackslash{}lemma\_CD\_chernoff\_adaptive.md}
-
\texttt{G:\textbackslash{}Xiaogan\_Supercomputing\_data\textbackslash{}SCX\textbackslash{}theory\textbackslash{}explorations\textbackslash{}lemma\_EF\_lowerbound\_aggregation.md}
-
\texttt{G:\textbackslash{}Xiaogan\_Supercomputing\_data\textbackslash{}SCX\textbackslash{}theory\textbackslash{}explorations\textbackslash{}cluster\_consistency\_v3.md}
-
\texttt{G:\textbackslash{}Xiaogan\_Supercomputing\_data\textbackslash{}SCX\textbackslash{}theory\textbackslash{}explorations\textbackslash{}deep\_math\_connections.md}

\begin{center}\rule{0.5\linewidth}{0.5pt}\end{center}

\subsection{1. Theorem 1 (Noise Detection
Guarantee)}\label{theorem-1-noise-detection-guarantee}

\subsubsection{Lemma 1 (Mean Separation)}\label{lemma-1-mean-separation}

\textbf{Claim}: E{[}C\textbar noise,x{]} = 1 -
E{[}C\textbar clean,x{]}/(K-1)

\textbf{Verification}: PASS.

The proof expands: - E{[}C\textbar noise,x{]} = (1/M) sum\_m P(f\_m(x)
!= y \textbar{} noise, x) - P(f\_m(x) != y \textbar{} noise, x) =
sum\_\{c != y\emph{\} P(y=c\textbar noise) } P(f\_m(x) != c \textbar{}
x) - = sum\_\{c != y\emph{\} (1/(K-1)) } (1 - P(f\_m(x)=c\textbar x)) -
= 1 - (1/(K-1)) * sum\_\{c != y\emph{\} P(f\_m(x)=c\textbar x) - = 1 -
(1/(K-1)) } (1 - P(f\_m(x)=y\emph{\textbar x)) - = 1 - (1/(K-1)) }
E{[}e\_m\textbar clean,x{]}

Averaging over m gives: E{[}C\textbar noise,x{]} = 1 -
E{[}C\textbar clean,x{]}/(K-1). The algebra is correct. Each step uses
only: (A4) uniform noise over K-1 classes, linearity of expectation, and
identity sum\_\{c != y*\} P(f\_m=c\textbar x) =
E{[}e\_m\textbar clean,x{]}.

The separation condition mu\_s \textless{} (K-1)/K is correctly derived
from: clean mean \textless= mu\_s, noise mean \textgreater= 1 -
mu\_s/(K-1), requiring mu\_s \textless{} 1 - mu\_s/(K-1) =\textgreater{}
mu\_s * K/(K-1) \textless{} 1 =\textgreater{} mu\_s \textless{} (K-1)/K.

The optimal threshold theta* = (1/2)(1 + mu\_s*(K-2)/(K-1)) solves theta
- mu\_s = 1 - mu\_s/(K-1) - theta correctly.

\subsubsection{Lemma 2 (FPR)}\label{lemma-2-fpr}

\textbf{Claim}: P(C \textgreater{} theta \textbar{} clean, s) \textless=
exp(-2M(theta - mu\_s)\^{}2)

\textbf{Verification}: PASS.

Standard Hoeffding inequality for bounded {[}0,1{]} independent r.v.s.
The conditioning chain: - Fix x in s: E{[}C\textbar clean,x{]}
\textless= mu\_s by (A5) - \{e\_m\} conditionally independent given x by
(A2), bounded by {[}0,1{]} by (A3) - P(C - E{[}C\textbar x{]}
\textgreater{} theta - E{[}C\textbar x{]} \textbar{} clean, x)
\textless= exp(-2M(theta - E{[}C\textbar x{]})\^{}2) \textless=
exp(-2M(theta - mu\_s)\^{}2) - The last step uses: theta \textgreater{}
mu\_s \textgreater= E{[}C\textbar x{]}, so (theta - E{[}C\textbar x{]})
\textgreater= (theta - mu\_s) - Marginalizing over x in s preserves the
bound via sup

All conditions for Hoeffding are satisfied.

\subsubsection{Lemma 3 (TPR)}\label{lemma-3-tpr}

\textbf{Claim}: P(C \textless= theta \textbar{} noise, s) \textless=
exp(-2M(1 - C\_bal*mu\_s/(K-1) - theta)\^{}2)

\textbf{Verification}: PASS, with careful reading.

The proof conditions on (x, c) where c is the noise label: - Given (x,
c), \{e\_m\} are conditionally independent (A1, A2) with e\_m
\textasciitilde{} Bern(1 - mu\_\{c,m\}(x)) - E{[}C\textbar x,c{]} = 1 -
mu\_c(x) \textgreater= 1 - C\_bal\emph{mu\_s/(K-1) by (A6) - The theorem
assumes theta \textless{} 1 - C\_bal}mu\_s/(K-1), so Hoeffding applies -
P(C \textless= theta \textbar{} x, c) \textless=
exp(-2M(E{[}C\textbar x,c{]} - theta)\^{}2) \textless= exp(-2M(1 -
C\_bal*mu\_s/(K-1) - theta)\^{}2) - Averaging over c (uniform over K-1
classes) preserves the bound

Key subtlety: the bound uses mu\_c(x) \textless=
C\_bal\emph{mu\_s/(K-1), so E{[}C\textbar x,c{]} \textgreater= 1 -
C\_bal}mu\_s/(K-1). The Hoeffding gap (E{[}C\textbar x,c{]} - theta) is
minimized (worst case) when E{[}C\textbar x,c{]} is smallest, which
occurs at the max mu\_c(x). The A6 assumption correctly captures this
via C\_bal.

\subsubsection{F1 bound}\label{f1-bound}

\textbf{Claim}: F1 \textgreater= 1 - (1/eta) sum\_s rho\_s exp(-2M
Delta\_s\^{}2)

\textbf{Verification}: PASS.

Algebraic derivation: - F1 = 2\emph{eta}TPR / (eta(1+TPR) + (1-eta)FPR)
{[}verified via Precision-Recall algebra{]} - Substituting TPR
\textgreater= 1-delta\_1, FPR \textless= delta\_2 gives denominator
\textgreater= eta - 1 - F1 \textless= delta\_1 + (1-eta)delta\_2/eta -
Since delta\_1, delta\_2 \textless= sum rho\_s exp(-2M Delta\_s\^{}2),
the final bound follows

The denominator is strictly positive: eta(2-FNR)+(1-eta)FPR
\textgreater= eta \textgreater{} 0 for eta \textgreater{} 0.

\begin{center}\rule{0.5\linewidth}{0.5pt}\end{center}

\subsection{2. Theorem 2 (Weak Features)}\label{theorem-2-weak-features}

\subsubsection{Fano inequality}\label{fano-inequality}

\textbf{Claim}: P(S\_hat != S) \textgreater= (H(S) - delta - log 2) /
log K

\textbf{Verification}: PASS.

Standard Fano inequality: P(S\_hat != S) \textgreater=
(H(S\textbar phi(X)) - log 2) / log \textbar S\textbar. Since I(phi;S)
\textless= delta by the delta-weak feature definition: H(S\textbar phi)
= H(S) - I(phi;S) \textgreater= H(S) - delta. Substituting gives:
P(S\_hat != S) \textgreater= (H(S) - delta - log 2) / log K.

The feature-dependent Markov chain Z -\textgreater{} S -\textgreater{} X
-\textgreater{} phi(X) correctly implies I(phi;Z) \textless= I(S;Z)
\textless= H(Z) by the data processing inequality.

\subsubsection{Pinsker inequality}\label{pinsker-inequality}

\textbf{Claim}: TV(P, P\_tilde) \textless= sqrt(delta/2)

\textbf{Verification}: PASS.

P\_tilde(phi, S) = P(phi) * P(S). KL(P \textbar\textbar{} P\_tilde) =
I(phi; S) = delta. Pinsker: TV \textless= sqrt(KL/2) = sqrt(delta/2).

\subsubsection{F1 bound}\label{f1-bound-1}

\textbf{Claim}: F1\_SCX \textless= F1\_base + C\_F * sqrt(delta/2)

\textbf{Verification}: PASS (with caveat on Lipschitz constant).

The reasoning chain: 1. Under P\_tilde, SCX detection score degrades to
loss baseline: NS(x) \textasciitilde{} max\_m D\_m(x) 2. Hence
F1\_P\_tilde(SCX) = F1\_base 3. \textbar F1\_P(SCX) -
F1\_P\_tilde(SCX)\textbar{} \textless= C\_F * TV(P\_pred,
P\_tilde\_pred) by Lipschitz property 4. TV(P\_pred, P\_tilde\_pred)
\textless= TV(P, P\_tilde) by data processing inequality 5. \textless=
sqrt(delta/2)

The Lipschitz analysis of F1(TP, FP, FN) = 2TP/(2TP+FP+FN) is partial
derivatives and the document acknowledges that C\_F depends on
precision/recall regime. The specific values (C\_F \textless= 2 for
precision, recall \textgreater= 0.1) are heuristic estimates, not
rigorous bounds. This is flagged as an approximate constant in the text.

\textbf{Minor issue}: The AUC bound's TV-to-conditional-TV step uses
division by eta and 1-eta: \textbar P(A\textbar Z=1) -
P\_tilde(A\textbar Z=1)\textbar{} \textless= TV(P, P\_tilde)/eta. This
requires eta \textgreater{} 0 (non-zero noise rate). For eta
-\textgreater{} 0 the bound diverges, which the document correctly
acknowledges.

\begin{center}\rule{0.5\linewidth}{0.5pt}\end{center}

\subsection{3. Theorem 3
(Unidentifiability)}\label{theorem-3-unidentifiability}

\subsubsection{K=2 construction}\label{k2-construction}

\textbf{Claim}: P\_noise(x, y, \{f\_m\}) = P\_hard(x, y, \{f\_m\})

\textbf{Verification}: PASS.

World A (noise): s1 has noise rate eta, clean expert error epsilon\_1.
s2 has no noise. World B (hard): s1 has random true label
(P(y\emph{=0)=1-eta, P(y}=1)=eta), experts biased toward class 0.

Equality verification: - P\_A(y=0\textbar s1) = (1-eta)\emph{1 + eta}0 =
1-eta - P\_B(y=0\textbar s1) = P(y\emph{=0\textbar s1) = 1-eta -
P\_A(f\_m=0\textbar s1) = 1-epsilon\_1 (expert accuracy against true
label y}=0) - P\_B(f\_m=0\textbar s1) = (1-eta)(1-epsilon\_1) +
eta(1-epsilon\_1) = 1-epsilon\_1

Joint factorization P(x,y,\{f\_m\}) = P(x) * P(y\textbar x) * prod\_m
P(f\_m\textbar x) holds by conditional independence in both
constructions. State s2 is identical in both worlds.

\subsubsection{K\textgreater2 construction}\label{k2-construction-1}

\textbf{Claim}: The construction works for all K \textgreater= 2.

\textbf{Verification}: PASS.

For K \textgreater{} 2, the construction uses completely random experts
(independent of y\emph{): - World A: y} = 0, noise uniform over K-1
classes. Expert predicts 0 with prob 1-epsilon\_1 - World B: y*
non-deterministic (same distribution as y in World A). Expert is pure
random

P\_A(y=0\textbar s1) = 1-eta = P\_B(y=0\textbar s1)
P\_A(f\_m=0\textbar s1) = 1-epsilon\_1 = P\_B(f\_m=0\textbar s1)
Independence of y and f\_m holds in both worlds by construction. Joint
distribution matches for all K \textgreater= 2.

\subsubsection{Error bound eta*rho/2}\label{error-bound-etarho2}

\textbf{Claim}: max(Error\_A, Error\_B) \textgreater= eta*rho/2

\textbf{Verification}: PASS.

Let a be the fraction of ambiguous samples (\{x in s1, y=1\}) flagged as
noise by the algorithm. Since observables are identical, a is the same
in both worlds. - World A: ambiguous samples are truly noise. Error
contribution \textgreater= (1-a) * eta * rho - World B: ambiguous
samples are truly clean. Error contribution \textgreater= a * eta * rho
- max(error\_A, error\_B) \textgreater= eta\emph{rho } max(1-a, a)
\textgreater= eta*rho/2

The bound is tight: achieved at a = 1/2 (random guess on ambiguous set).

\begin{center}\rule{0.5\linewidth}{0.5pt}\end{center}

\subsection{4. Theorem 4' (Exact Constant
Minimax)}\label{theorem-4-exact-constant-minimax}

\subsubsection{Bahadur-Rao saddlepoint}\label{bahadur-rao-saddlepoint}

\textbf{Claim}: lambda* = log(theta(1-p)/(p(1-theta))),
sigma\^{}2(theta) = theta(1-theta)

\textbf{Verification}: PASS.

For Bern(p): psi(lambda) = log(1-p + p\emph{e\^{}lambda) psi'(lambda) =
p}e\textsuperscript{lambda/(1-p+p\emph{e\^{}lambda). Setting = theta:
p}e}lambda = theta(1-p+p\emph{e\^{}lambda) =\textgreater{} e\^{}lambda =
theta(1-p)/(p(1-theta)) =\textgreater{} lambda} =
log(theta(1-p)/(p(1-theta)))

Tilted variance: psi'\,'(lambda\emph{) = theta(1-theta). Verified by
direct substitution: 1-p+p}e\^{}(lambda*) = (1-p)/(1-theta), giving
psi'\,' = theta(1-theta).

\subsubsection{Chernoff information}\label{chernoff-information}

\textbf{Claim}: kappa = KL(theta\emph{\textbar\textbar p0) =
KL(theta}\textbar\textbar p1) where theta* =
log((1-p0)/(1-p1))/log(p1(1-p0)/(p0(1-p1)))

\textbf{Verification}: PASS.

The equality KL(theta\textbar\textbar p0) = KL(theta\textbar\textbar p1)
gives: theta\emph{log(p1/p0) = (1-theta)}log((1-p0)/(1-p1)) Solving:
theta* = log((1-p0)/(1-p1)) / log(p1(1-p0)/(p0(1-p1)))

Verified theta* is in (p0, p1) by strict convexity of KL in its first
argument and sign change at endpoints.

\subsubsection{Adaptive threshold}\label{adaptive-threshold}

\textbf{Claim}: theta\_opt = theta* + (1/M)\emph{log((1-eta)/eta)/D} +
O(1/M\^{}2)

\textbf{Verification}: PASS.

Minimizing (1/2)FNR(theta) + (1-eta)/(2eta)FPR(theta) gives first-order
condition: KL(theta\textbar\textbar p0) - KL(theta\textbar\textbar p1) =
(1/M)*log((1-eta)/eta)

Key observation: KL'\,`(theta\textbar\textbar p) = 1/(theta(1-theta)) is
INDEPENDENT of p.~Therefore: KL(theta\emph{+delta\textbar\textbar p0) -
KL(theta}+delta\textbar\textbar p1) =
delta\emph{(KL'(theta}\textbar\textbar p0) -
KL'(theta\emph{\textbar\textbar p1)) + O(delta\^{}3) = delta}(lambda0* -
lambda1\emph{) + O(delta\^{}3) = delta}D* + O(delta\^{}3)

The delta\^{}2 terms cancel exactly. This is a clean mathematical fact
correctly exploited.

Therefore: delta = (1/M)\emph{log((1-eta)/eta)/D} + O(1/M\^{}2).

\subsubsection{O(1) prefactor}\label{o1-prefactor}

\textbf{Claim}: ((1-eta)/eta)\^{}s appears in both FPR and FNR
contributions.

\textbf{Verification}: PASS.

At theta\_opt = theta* + delta:
exp(-M\emph{KL(theta\_opt\textbar\textbar p0)) = e\^{}\{-M}kappa\} *
((1-eta)/eta)\^{}\{-(1-s)\}
exp(-M\emph{KL(theta\_opt\textbar\textbar p1)) = e\^{}\{-M}kappa\} *
((1-eta)/eta)\^{}s

where s = \textbar lambda1\emph{\textbar/D}, 1-s = lambda0\emph{/D}.

FPR term:
((1-eta)/(2eta))\emph{exp(-M}kappa)\emph{((1-eta)/eta)\^{}\{-(1-s)\}/(lambda0}sqrt(\ldots))
= exp(-M\emph{kappa)}((1-eta)/eta)\^{}s/(2\emph{lambda0}sqrt(\ldots))

FNR term:
(1/2)\emph{exp(-M}kappa)\emph{((1-eta)/eta)\^{}s/(\textbar lambda1}\textbar*sqrt(\ldots))

Both carry the identical factor ((1-eta)/eta)\^{}s. This cancellation is
verified algebraically.

\subsubsection{C\_min formula}\label{c_min-formula}

\textbf{Claim}: C\_min =
(eta/2)\emph{((1-eta)/eta)\^{}s}(1/lambda0\emph{+1/\textbar lambda1}\textbar)/sqrt(theta\emph{(1-theta}))

\textbf{Verification}: PASS.

From the weighted risk at the Bayes optimal threshold: R\_M =
w0\emph{alpha\_M + w1}beta\_M =
e\textsuperscript{\{-M\emph{kappa\}}((1-eta)/eta)}s\emph{(1/lambda0}+1/\textbar lambda1\emph{\textbar)/(2}sqrt(2pi\emph{M}theta\emph{(1-theta})))*(1+o(1))

lim e\^{}\{M\emph{kappa\}}sqrt(2pi\emph{M)}R\_M =
((1-eta)/eta)\^{}s\emph{(1/lambda0}+1/\textbar lambda1\emph{\textbar)/(2}sqrt(theta\emph{(1-theta})))
= C\_min/eta

Solving: C\_min = eta * {[}above expression{]} =
(eta/2)\emph{((1-eta)/eta)\^{}s}(1/lambda0\emph{+1/\textbar lambda1}\textbar)/sqrt(theta\emph{(1-theta}))

\subsubsection{Constant matching}\label{constant-matching}

\textbf{Claim}: Adaptive SCX constant = C\_min/eta

\textbf{Verification}: PASS.

From Lemma D section D.5: 1-F1(theta\_opt) \textasciitilde{}
e\textsuperscript{\{-M\emph{kappa\}}((1-eta)/eta)}s\emph{(1/lambda0}+1/\textbar lambda1\emph{\textbar)/(2}sqrt(2pi\emph{M}theta\emph{(1-theta})))

This is exactly C\_min/eta. The lower bound (Lemma E) proves no
algorithm can have a smaller constant. SCX with adaptive threshold
achieves it. Match is exact.

\begin{center}\rule{0.5\linewidth}{0.5pt}\end{center}

\subsection{5. Theorem 5 (Cluster
Consistency)}\label{theorem-5-cluster-consistency}

\subsubsection{All exponents are
negative}\label{all-exponents-are-negative}

\textbf{Claim}:
-c1\emph{n\_min}Delta\_min\textsuperscript{2/(sigma}2*d\_phi) is always
negative.

\textbf{Verification}: PASS.

The exponent contains: c1 \textgreater{} 0 (universal constant), n\_min
\textgreater{} 0 (minimum per-state sample size), Delta\_min
\textgreater{} 0 (minimum separation, fixed), sigma\^{}2 \textless{} inf
(sub-Gaussian proxy), d\_phi \textless{} inf (feature dimension). All
factors are positive, so the exponent is strictly negative. The bound
converges to 0 as n\_min -\textgreater{} inf.

Additionally verified: - Lemma S1:
P(\textbar\textbar epsilon\textbar\textbar\_2 \textgreater=
C1\emph{sigma}(sqrt(d\_phi)+sqrt(t))) \textless= 2\emph{exp(-t).
Exponent -t is always negative for t \textgreater{} 0. - Lemma S2:
Sub-Gaussian tail bound with negative exponent. - Lemma 2 peeling sum:
each term exp(-c}4\^{}j*t\^{}2) with negative exponent, summed
exponentially.

\subsubsection{All inequality
directions}\label{all-inequality-directions}

\textbf{Claim}: All inequalities point in correct directions.

\textbf{Verification}: PASS.

Verification table from Section 7.3 checked: - Lemma 1 Claim: P(event)
\textless= exp(-Delta\_min\textsuperscript{2/(8\emph{sigma\^{}2)). The
derivation: condition \textbar\textbar mu\_j-mu\_k +
epsilon\textbar\textbar{} \textless=
\textbar\textbar epsilon\textbar\textbar{} implies
2}epsilon}T(mu\_j-mu\_k) \textless=
-\textbar\textbar mu\_j-mu\_k\textbar\textbar\^{}2. Sub-Gaussian tail
gives the exp bound. Direction: upper bound (correct for tail). - Lemma
3: \textbar\textbar phi - theta\_\{j\_k\}\textbar\textbar{} \textless=
Delta\_min/2 \textless= \textbar\textbar phi -
theta\_\{j'\}\textbar\textbar. Verifies 1/8 + 3/8 = 1/2 and 7/8 - 3/8 =
1/2. Both sides with correct inequality. - Lemma 2 (3):
lambda\emph{\textbar\textbar theta\_hat - theta}\textbar\textbar\^{}2
\textless= \textbar(Pn-P)(f\_theta\_hat - f\_theta\emph{)\textbar. From
W\_n(theta\_hat) \textless= W\_n(theta}), expanding and using quadratic
lower bound. Direction: correct. - Lemma 4: P(no restart lands in B)
\textless= (1-p0)\^{}R \textless= n\^{}\{-c\}. Direction: correct.

\subsubsection{NP-hard gap honestly
discussed}\label{np-hard-gap-honestly-discussed}

\textbf{Claim}: The NP-hard gap of k-means is honestly acknowledged and
addressed.

\textbf{Verification}: PASS.

Section 6 and 9.2 explicitly: 1. State that k-means is NP-hard in worst
case 2. Lemma 4 proves Lloyd's with R = O(log n) random restarts finds
global minimizer under strong separation 3. Explicitly state the
consequence if strong separation fails: ``if the strong separation
condition does not hold\ldots{} the theorem's guarantee degrades'' 4.
Cite relevant literature (Ostrovsky et al.~2013, Kumar \& Kannan 2010)
5. Do not claim a guarantee for the weak separation regime

This is an honest, well-structured treatment of the computational
hardness.

\begin{center}\rule{0.5\linewidth}{0.5pt}\end{center}

\subsection{6. Numerical Verification}\label{numerical-verification}

\subsubsection{Test 1: p0=0.10, p1=0.60,
eta=0.10}\label{test-1-p00.10-p10.60-eta0.10}

\begin{longtable}[]{@{}lll@{}}
\toprule\noalign{}
Quantity & Computed Value & Verification \\
\midrule\noalign{}
\endhead
\bottomrule\noalign{}
\endlastfoot
theta* & 0.31158 & Closed-form formula \\
kappa & 0.1696 & Matches Lemma C table (0.1696) \\
lambda0* & 1.4040 & \textgreater{} 0, correct direction \\
& lambda1* & \\
D* & 2.6021 & lambda0* + \\
s & 0.4604 & \\
C\_min & 0.4573 & Formula evaluated \\
C\_min/eta & 4.573 & \\
C\_SCX & 4.573 & C\_SCX = C\_min/eta: MATCH \\
\end{longtable}

kappa/2Delta\^{}2 comparison: 2\emph{(p1-p0)\^{}2 = 2}0.25 = 0.50.
kappa/0.50 = 0.1696/0.50 = 0.3392, matching Lemma C table.

\subsubsection{Test 2: p0=0.05, p1=0.80,
eta=0.05}\label{test-2-p00.05-p10.80-eta0.05}

\begin{longtable}[]{@{}lll@{}}
\toprule\noalign{}
Quantity & Computed Value & Verification \\
\midrule\noalign{}
\endhead
\bottomrule\noalign{}
\endlastfoot
theta* & 0.35979 & Closed-form formula \\
kappa & 0.4576 & Matches Lemma C table (0.4574, minor rounding) \\
lambda0* & 2.368 & \textgreater{} 0 \\
& lambda1* & \\
D* & 4.330 & lambda0* + \\
s & 0.4531 & \\
C\_min & 0.1863 & Formula evaluated \\
C\_min/eta & 3.727 & \\
C\_SCX & 3.727 & C\_SCX = C\_min/eta: MATCH \\
\end{longtable}

\subsubsection{Test 3: p0=0.20, p1=0.50,
eta=0.30}\label{test-3-p00.20-p10.50-eta0.30}

\begin{longtable}[]{@{}lll@{}}
\toprule\noalign{}
Quantity & Computed Value & Verification \\
\midrule\noalign{}
\endhead
\bottomrule\noalign{}
\endlastfoot
theta* & 0.33904 & Closed-form formula \\
kappa & 0.0527 & Matches Lemma C table (0.0528, minor rounding) \\
lambda0* & 0.7186 & \textgreater{} 0 \\
& lambda1* & \\
D* & 1.3859 & lambda0* + \\
s & 0.4815 & \\
C\_min & 1.373 & Formula evaluated \\
C\_min/eta & 4.578 & \\
C\_SCX & 4.578 & C\_SCX = C\_min/eta: MATCH \\
\end{longtable}

All three numerical test cases confirm the identity C\_SCX = C\_min/eta.

\begin{center}\rule{0.5\linewidth}{0.5pt}\end{center}

\subsection{7. Minor Observations and
Notes}\label{minor-observations-and-notes}

\begin{enumerate}
\def\labelenumi{\arabic{enumi}.}
\item
  \textbf{Lemma A lattice correction}: The Bahadur-Rao expansion for
  Bernoulli distributions includes a lattice correction factor (1 -
  e\^{}\{-lambda\emph{\})\^{}\{-1\} not present in the continuous case.
  The document correctly notes both the 1/lambda} form and the lattice
  correction form, showing they are asymptotically equivalent. When used
  consistently in FPR/FNR ratios (which share the same theta*), the
  correction largely cancels. This is mathematically sound.
\item
  \textbf{Theorem 2 Lipschitz constant C\_F}: The Lipschitz constant for
  F1 as a function of (TP, FP, FN) depends on the operating regime. The
  document gives heuristic values (C\_F \textless= 3 for precision,
  recall \textgreater= 0.1; C\_F \textless= 1 for precision, recall
  \textgreater= 0.5). These are not rigorous bounds derived in the
  document, but the claim that a Lipschitz constant exists is correct.
  The specific values are practical estimates.
\item
  \textbf{Theorem 4' non-asymptotic bounds}: Section 6(d) claims
  finite-M bounds with explicit constants K1, K2, but states these come
  from Lemma A (Stirling error) and Lemma B (quadratic remainder). The
  lemmas provide error terms of the form O(1/M), but the explicit K1, K2
  constants are referenced rather than fully computed. The asymptotic
  limit (M -\textgreater{} inf) is fully rigorous.
\item
  \textbf{Theorem 5 irreducible error}: The document acknowledges that
  even with perfect center estimates, points with
  \textbar\textbar epsilon\textbar\textbar{} \textgreater=
  3*Delta\_min/8 can be misclassified. This is correctly identified as a
  data property, not an estimator failure.
\end{enumerate}

\begin{center}\rule{0.5\linewidth}{0.5pt}\end{center}

\subsection{8. Overall Mathematical
Verdict}\label{overall-mathematical-verdict}

All mathematical derivations across Theorems 1-5 and supporting lemmas
are \textbf{CORRECT}.

\begin{longtable}[]{@{}
  >{\raggedright\arraybackslash}p{(\linewidth - 4\tabcolsep) * \real{0.2903}}
  >{\raggedright\arraybackslash}p{(\linewidth - 4\tabcolsep) * \real{0.2903}}
  >{\raggedright\arraybackslash}p{(\linewidth - 4\tabcolsep) * \real{0.4194}}@{}}
\toprule\noalign{}
\begin{minipage}[b]{\linewidth}\raggedright
Theorem
\end{minipage} & \begin{minipage}[b]{\linewidth}\raggedright
Verdict
\end{minipage} & \begin{minipage}[b]{\linewidth}\raggedright
Key Findings
\end{minipage} \\
\midrule\noalign{}
\endhead
\bottomrule\noalign{}
\endlastfoot
Theorem 1 (Noise Detection) & PASS & All lemmas algebraically verified;
Hoeffding conditions satisfied; F1 bound derivations correct \\
Theorem 2 (Weak Features) & PASS & Fano/Pinsker/DPI applied correctly;
TV-to-F1 Lipschitz argument valid (constants are heuristic but not
essential) \\
Theorem 3 (Unidentifiability) & PASS & K=2 and K\textgreater2
constructions both verified; joint distribution equality holds; error
bound eta*rho/2 is tight \\
Theorem 4' (Exact Constant) & PASS & Bahadur-Rao verified in full;
Chernoff theta* closed-form correct; adaptive threshold derivation
exploits KL'\,' independence of p; C\_min formula matches all three
numerical tests; C\_SCX = C\_min/eta verified \\
Theorem 5 (Cluster Consistency) & PASS & All exponents negative; all
inequality directions verified; NP-hard gap honestly discussed and
addressed under strong separation \\
Lemma A (Bahadur-Rao) & PASS & Full saddlepoint computation verified;
tilted distribution correctly identified as Bern(theta); lattice
correction discussed \\
Lemma B (F1 expansion) & PASS & F1 formula algebra verified; denominator
positivity checked; expansion into FNR/2 + (1-eta)FPR/(2eta) correct;
remainder bound derived \\
Lemma C (Chernoff info) & PASS & Theta* closed-form derived correctly;
KL'\,' independence noted; numerical table entries verified against
formulas \\
Lemma D (Adaptive threshold) & PASS & First-order condition derived
correctly; O(1/M) shift computed; (1-eta)/eta)\^{}s cancellation
verified \\
Lemma E (Lower bound) & PASS & Bayes test reduction via Neyman-Pearson
is correct; Bahadur-Rao applied; C\_min formula derived explicitly and
matches achievable constant \\
Lemma F (Multi-state) & PASS & Additivity by law of total expectation;
bottleneck state dominates asymptotic; correct handling of kappa\_min
states \\
\end{longtable}

\end{document}
